This work is devoted to establishing a rigorous lower bound for the Chv\'atal--Sankoff constant $\gamma_\sigma$ --- the limit of the ratio of the expected length of the longest common subsequence (LCS) of two i.i.d.\ uniformly random strings of length $n$ over an alphabet of size $\sigma$ to $n$ as $n \to \infty$. The value of $\gamma_\sigma$ is known only as an interval: for $\sigma = 2$ it is proven that $\gamma_2 \in [0.7927;\, 0.8263]$ (Heineman et al., 2024; Lueker, 2009), while numerical estimates suggest $\gamma_2 \approx 0.811$.

A new method for computing lower bounds is proposed --- windowed dynamic programming with bounded look-ahead $N$ and maximal shift $\mathrm{MAX\_SH}$. Proof idea: the DP computes the optimal mean number of matches achievable by a bounded-look-ahead strategy over $K$ steps; by induction and Fekete's superadditivity lemma, the normalized limit $f_K(0,0,0)/K$ exists and is bounded above by $\gamma_\sigma$, since every such strategy is realizable on actual random strings. Algorithm complexity: $O(K \cdot \mathrm{MAX\_SH} \cdot \sigma^{2N+2})$ time, $O(\mathrm{MAX\_SH} \cdot \sigma^{2N})$ memory. The method is universally applicable for any $\sigma \geq 2$. A C++ implementation using AVX2 and single-precision arithmetic allows $N = 12$, $\mathrm{MAX\_SH} = 30$, $K = 40\,000$ with $\approx 2$~GB memory and $\approx 5$~hours on a single core.

New record lower bounds are obtained: $\gamma_2 \geq 0.7964$ (improvement of $+0.48\%$ over Heineman 2024), as well as records for all $\sigma \in \{3, \ldots, 10\}$ with improvements from $+0.17\%$ to $+5.42\%$ compared to the best known results. All bounds are independently verified by Monte-Carlo simulation on strings of length up to $20\,000$ (with $500$ trials per $\sigma$).
